% =====================================================================
%  Цель и задание лабораторной работы.
% =====================================================================

\textbf{Цель:} получить практические навыки анализа защищённости
реального удалённого сервера, доступного из сети Internet по SSH и
HTTP(S), на основании данных, накапливаемых системами журналирования
в течение нескольких дней естественной эксплуатации.

\vspace{1em}

\textbf{Задание}:
\begin{enumerate}
  \item Установить на удалённом сервере веб-сервер и разместить на нём
  тестовую страницу, доступную из сети Internet.
  \item Настроить систему журналирования на веб-сервере.
  \item Выполнить анализ данных логов различных подсистем, в том числе
  SSH и веб-сервера удалённого сервера, за несколько дней.
  \item Оформить отчёт о проделанной работе.
\end{enumerate}

\textbf{Постановка задачи.} В качестве удалённого сервера используется
виртуальный сервер (VPS) под управлением Ubuntu~24.04 с публичным
IP-адресом и доменным именем \texttt{test-lab4.work.gd}. В роли
веб-сервера выбран \texttt{Caddy} -- он раздаёт тестовую страницу
(одностраничное приложение на React) по протоколу HTTPS с автоматически
выпущенным сертификатом Let's Encrypt. Сервер намеренно оставлен
доступным из интернета без дополнительного ограничения (кроме
управления по ключу SSH вместо пароля) на несколько дней, чтобы в логах
SSH-демона и веб-сервера накопился настоящий фоновый интернет-трафик
(автоматические сканеры, попытки подбора пароля и т.\,п.), а не
искусственно смоделированные обращения.
